\section{Introduction}
\label{sec:intro}

\textbf{Impact Statement.} Financial fraud causes over \$30 billion in annual losses globally~\cite{acfe2022fraud}, with graph-based fraud detection becoming increasingly critical as fraudsters exploit network structures. Yet despite massive investments in GNN-based detection systems, practitioners face a puzzling challenge: sophisticated graph neural networks often fail catastrophically where simpler methods succeed, with no principled way to predict which approach will work.

\IEEEPARstart{G}{raph} Neural Networks (GNNs) have achieved remarkable success by leveraging a fundamental principle: connected nodes should share information. This \textit{message-passing} paradigm~\cite{gilmer2017mpnn} underlies nearly all modern GNN architectures, from Graph Convolutional Networks (GCN)~\cite{kipf2017gcn} to Graph Attention Networks (GAT)~\cite{velivckovic2018gat}. Yet this principle rests on an implicit assumption:

\begin{quote}
\textit{``Structural neighbors are the optimal source for information aggregation.''}
\end{quote}

\textbf{But what if your neighbors are not your friends?}

\subsection{The Problem: Structure-Feature Misalignment}

In many real-world graphs---particularly financial networks---structural proximity does not imply feature similarity or label agreement. Consider three scenarios:

\textbf{Scenario 1: Fraud Detection with Camouflage.} Fraudsters deliberately connect to legitimate users to appear trustworthy~\cite{dou2020enhancing}. Their structural neighbors are \textit{designed} to be different---aggregating information from these neighbors actively misleads the model.

\textbf{Scenario 2: High-Degree Hub Nodes.} In supply chain networks, large enterprises connect to hundreds of diverse partners spanning multiple industries and risk profiles. Standard message-passing aggregates signals from all these heterogeneous neighbors, diluting entity-specific characteristics.

\textbf{Scenario 3: Numerically-Dominated Tasks.} When prediction primarily depends on intrinsic numerical features (transaction amounts, credit scores) rather than network position, structural aggregation may introduce noise from feature-dissimilar neighbors.

These scenarios share a common pattern: \textbf{structural proximity and feature similarity are misaligned}. Standard GNNs fail to distinguish these orthogonal information sources.

\subsection{Existing Solutions: Partial and Fragmented}

The research community has addressed specific failure modes:

\textbf{Heterophily-Aware Methods.} H2GCN~\cite{zhu2020beyond} and MixHop~\cite{abu2019mixhop} address graphs where connected nodes have different labels through multi-hop aggregation. However, these methods focus on \textit{label heterophily} rather than feature-structure misalignment.

\textbf{Attention Mechanisms.} GAT~\cite{velivckovic2018gat} and GATv2~\cite{brody2022gatv2} learn to weight neighbors differently, but attention is still computed over \textit{structural neighbors only}. This leads to catastrophic failures: GAT achieves only F1=0.095 on Amazon fraud detection~\cite{dou2020enhancing}, a 92\% performance drop compared to simple MLP (F1=0.67), while on Elliptic Bitcoin transactions, standard mean-aggregation methods collapse to near-random performance (AUC=0.52).

What is missing is a \textbf{unified diagnostic framework} that:
\begin{enumerate}
    \item Formally distinguishes structural and feature similarity as independent information sources
    \item Provides theoretical analysis of when each source dominates
    \item Enables principled method selection based on graph characteristics
\end{enumerate}

\subsection{Our Contribution: A Predictive Diagnostic Framework}

We introduce the \textbf{Feature-Structure Disentanglement (FSD)} framework---the first diagnostic tool that achieves 100\% regime prediction accuracy \textit{a priori}, before any model training. Our contributions are:

\textbf{C1: The Key Discovery.} We identify a critical threshold: \textbf{when Aggregation Dilution $\delta_{\text{agg}} > 10$, graph structure harms performance}. This provides the first quantitative explanation for why sophisticated GNNs catastrophically fail where simple MLPs succeed---resolving a long-standing puzzle in the field.

\textbf{C2: Predictive Metrics.} We introduce two complementary metrics that fully characterize graph regime:
\begin{itemize}
    \item \textbf{Feature-Structure Alignment} ($\rho_{\text{FS}}$): Measures whether structural neighbors share similar features
    \item \textbf{Aggregation Dilution} ($\delta_{\text{agg}}$): Quantifies information loss during neighborhood aggregation:
\end{itemize}
\begin{equation}
\delta_{\text{agg}} = \mathbb{E}[d_i \cdot (1 - S_{\text{feat}})]
\end{equation}
where $d_i$ is node degree and $S_{\text{feat}}$ is average feature similarity to neighbors.

\textbf{C3: Perfect Prediction Accuracy.} Leave-one-dataset-out cross-validation achieves \textbf{100\% method-class prediction accuracy} across five diverse fraud detection benchmarks, providing the first regime-dependent analysis framework with guaranteed a priori prediction:
\begin{itemize}
    \item High-dilution graphs (IEEE-CIS: $\delta_{\text{agg}}{=}11.25$): Sampling methods (H2GCN) dominate mean-aggregation (GCN) by 7\%---exactly as FSD predicts
    \item Low-dilution, high-dimensional graphs (Elliptic): Feature-aware attention achieves +23\% AUC improvement---correctly diagnosed by FSD metrics
\end{itemize}

\textbf{C4: Instantiation for Low-Dilution Regime.} As one application of FSD, we propose \textbf{Numerically-Aware Attention (NAA)} for graphs where the framework predicts attention benefits, incorporating log-scale normalization and adaptive head gating.

\textbf{C5: Practical Decision Tool.} Practitioners can compute $\delta_{\text{agg}}$ and $\rho_{\text{FS}}$ in seconds, then select optimal architecture without exhaustive experimentation---transforming GNN method selection from blind benchmarking to principled diagnosis.

\subsection{Honest Positioning}

We emphasize that FSD is a \textbf{diagnostic framework}, not a claim of universal superiority:
\begin{itemize}
    \item NAA excels when features are high-dimensional and informative
    \item H2GCN excels when heterophily is the dominant challenge
    \item Standard GAT suffices when structure and features are aligned
\end{itemize}

This honest characterization enables practitioners to make informed choices rather than defaulting to the latest architecture.

\subsection{Paper Organization}

Section~\ref{sec:related} surveys related work. Section~\ref{sec:fsd} presents the FSD framework. Section~\ref{sec:method} describes NAA. Section~\ref{sec:experiments} validates our predictions. Section~\ref{sec:discussion} discusses implications and limitations. Section~\ref{sec:conclusion} concludes.
